\section{A vanishing radial margin}
\label{sec:angular-cubic}

For a fixed positive margin, the normalized area coefficient is
$2\eta-1/240$. We now let the margin tend to zero with the tilt and
obtain the coefficient $-1/240$ in a single family of finite triangular
covers. The extra radial enlargement is cubic in the tilt. Its positive
cubic area term explains why coverage persists on a much wider interval
than the interval on which we prove a saving.

The fixed-margin theorem, Theorem~\ref{ext:general-cap}, cannot be
applied uniformly when $\eta$ tends to zero. Instead we use the explicit
coverage budget in Corollary~\ref{cor:quantitative-coverage} and the
uniform area estimate in Corollary~\ref{cor:variable-margin-area}.
We keep the edge-two tetrahedron $K$ and the same slope coordinates.

\begin{proposition}\label{prop:angular-cubic}
For every $0<\tau\le1$, the tetrahedron $K$ admits a cover by
$2\lceil8/\tau^2\rceil$ cylinders with compact, nondegenerate triangular
perpendicular bases. Their total base area $C_\tau$ satisfies
\begin{equation}\label{eq:angular-cubic-cost}
 \left|\frac{C_\tau}{H}-\frac12+\frac{\tau^2}{240}
                   -\frac{25}{2}\tau^3\right|
 \le2\tau^4\qquad(0<\tau\le1/50).
\end{equation}
In particular, $C_\tau<A_{\min}(K)/2$ for $0<\tau\le1/4000$, and
\[
 \frac{C_\tau}{A_{\min}(K)}
 =\frac12-\frac{\tau^2}{240}+O(\tau^3)
 \qquad(\tau\downarrow0).
\]
\end{proposition}
\begin{proof}
Put
\[
 N=\left\lceil\frac8{\tau^2}\right\rceil,\qquad
 \Delta=\frac2N\le\frac{\tau^2}{4},\qquad
 q_j=-1+j\Delta\quad(0\le j\le N).
\]
Use the secant coefficients \eqref{eq:axis-coefficients}, and set
\begin{equation}\label{eq:cubic-caps}
 d_j=\max_{q_j\le q\le q_{j+1}}d(q),\qquad
 T_j=\frac12+\tau^2d_j+\frac{25}{4}\tau^3
 \qquad(0\le j<N).
\end{equation}
The intercept triangles and directions are
\eqref{eq:first-family}--\eqref{eq:second-family}, with $\eps$ replaced
by $\tau$ and these cutoffs. They give exactly $2N$ cylinders.
Their intercept triangles have positive area, and the normal component
of each axis is one, so their perpendicular bases are compact
nondegenerate triangles as before.

\emph{Coverage.} For $0<\tau\le1/2$, take $\eta=25\tau/4$ in
Corollary~\ref{cor:quantitative-coverage}. Its sufficient budget is
\[
 2\eta-(2+2\eta)\tau-\frac{\tau^2}{4}
 =\tau\left(\frac{21}{2}-\frac{51}{4}\tau\right)
 \ge\frac{33}{8}\tau>0.
\]
Thus the cylinders cover $K$ throughout this range.

For $1/2\le\tau\le1$, the enlarged triangles are already large enough
for the first family alone. Indeed, for $(x,y,Ht)\in K$ the first-crossing
argument selects a first-family sector with intercept height
\[
 0\le t_0=t-\tau x\beta_j
 \le t+\frac{\tau}{2}|x|
 \le t+\frac{\tau}{2}(1-t)\le1.
\]
Its cutoff obeys
$T_j\ge1/2+(25/4)\tau^3\ge41/32>1$.
The selected cylinder therefore contains the point. Both arguments use
non-strict angular membership conditions, so they cover the closed
body, including its boundary.

\emph{Area.} For $0<\tau\le1/50$, the margin $\eta=25\tau/4$ lies
in $[0,1/8]$. Corollary~\ref{cor:variable-margin-area} consequently gives
\[
 \left|\frac{C_\tau}{H}-\frac12-
       \left(\frac{25}{2}\tau-\frac1{240}\right)\tau^2\right|
 \le2\tau^4,
\]
which is \eqref{eq:angular-cubic-cost}. This application is uniform
although the margin varies with the tilt. If $0<\tau\le1/4000$, then
\[
 \frac{25}{2}\tau+2\tau^2
 \le\frac1{320}+\frac1{8{,}000{,}000}<\frac1{240}.
\]
The negative quadratic term therefore exceeds both the positive cubic
term and the error, giving $C_\tau<H/2=A_{\min}(K)/2$.
\end{proof}

The equivalent edge-four construction, with its full coverage interval,
exact cylinder count and scaled area estimates, is recorded in
Remark~\ref{rem:angular-cubic-scale} of Appendix~\ref{ext:scaling-specializations}.
